\documentclass[11pt]{article}

\usepackage[margin=0.45in]{geometry}
\usepackage{amsmath,amssymb}
\usepackage{booktabs}
\usepackage{array}
\usepackage{xcolor}
\usepackage{tcolorbox}
\usepackage{tabularx}
\usepackage{makecell}
\usepackage{hyperref}
\usepackage{fancyhdr}
\usepackage{lastpage}

\definecolor{main}{HTML}{1F4E79}
\definecolor{softblue}{HTML}{EEF6FF}

\hypersetup{
    colorlinks=true,
    linkcolor=main,
    urlcolor=main
}

\newcolumntype{L}[1]{>{\raggedright\arraybackslash}p{#1}}
\newcolumntype{C}[1]{>{\centering\arraybackslash}p{#1}}
\newcolumntype{Y}{>{\centering\arraybackslash}X}

\newtcolorbox{distbox}[1]{
    colback=softblue,
    colframe=main,
    title=\textbf{#1},
    fonttitle=\large,
    arc=3mm,
    boxrule=0.8pt,
    left=3mm,
    right=3mm,
    top=2mm,
    bottom=1mm
}

% -------------------------------------------------
% Header / Footer Style
% -------------------------------------------------

\pagestyle{fancy}
\fancyhf{}

\setlength{\headheight}{18pt}
\setlength{\headsep}{18pt}

\lhead{\small University of Washington}
\chead{\small Department of Mathematics}
\rhead{\small MATH 394: Probability I}

\lfoot{\small Summer 2026}
\cfoot{\small Arman Jahangiri}
\rfoot{\small \thepage\ /\ \pageref{LastPage}}

\renewcommand{\headrulewidth}{0.4pt}
\renewcommand{\footrulewidth}{0.4pt}

\begin{document}

\begin{center}
{\Large\bfseries Probability Distribution Reference Sheet}\\[-1mm]
{\normalsize MATH 394: Probability I (Summer 2026)}\\[-1mm]
{\small Instructor: Arman Jahangiri, Department of Mathematics @ University of Washington}
\end{center}

\vspace{-2mm}

\begin{distbox}{Discrete Distributions I}
\scriptsize
\renewcommand{\arraystretch}{1.55}
\setlength{\tabcolsep}{2.2pt}
\begin{tabularx}{\textwidth}{@{}L{2.0cm}YYYYY@{}}
\toprule
\textbf{Quantity} & \textbf{Bernoulli} & \textbf{Binomial} & \textbf{Poisson} & \textbf{Geometric} & \textbf{Geometric, zero-based} \\
\midrule
Parameters
& $0\le p\le 1$
& $n\ge 1,\ 0\le p\le 1$
& $\lambda>0$
& $0<p\le 1$
& $0<p\le 1$ \\
PMF
& $P(X=k)=p^k(1-p)^{1-k}$
& $P(X=k)=\binom{n}{k}p^k(1-p)^{n-k}$
& $P(X=k)=e^{-\lambda}\dfrac{\lambda^k}{k!}$
& $P(X=k)=p(1-p)^{k-1}$
& $P(X=k)=p(1-p)^k$ \\
Support
& $k=0,1$
& $k=0,\dots,n$
& $k=0,1,2,\dots$
& $k=1,2,\dots$
& $k=0,1,2,\dots$ \\
Mean
& $p$
& $np$
& $\lambda$
& $\dfrac{1}{p}$
& $\dfrac{1-p}{p}$ \\
Variance
& $p(1-p)$
& $np(1-p)$
& $\lambda$
& $\dfrac{1-p}{p^2}$
& $\dfrac{1-p}{p^2}$ \\
MGF
& $1-p+pe^t$
& $(1-p+pe^t)^n$
& $\exp\{\lambda(e^t-1)\}$
& $\dfrac{pe^t}{1-(1-p)e^t}$
& $\dfrac{p}{1-(1-p)e^t}$ \\
\bottomrule
\end{tabularx}
\end{distbox}

\vspace{2mm}

\begin{distbox}{Discrete Distributions II}
\scriptsize
\renewcommand{\arraystretch}{1.55}
\setlength{\tabcolsep}{2.2pt}
\begin{tabularx}{\textwidth}{@{}L{2.0cm}YYYY@{}}
\toprule
\textbf{Quantity} & \textbf{Negative Binomial} & \textbf{Negative Binomial, failures before $r$ successes} & \textbf{Hypergeometric} & \textbf{Negative Hypergeometric} \\
\midrule
Parameters
& $r\ge 1,\ 0<p\le 1$
& $r\ge 1,\ 0<p\le 1$
& $w,b,n$
& $w,b,r$ \\
PMF
& $P(X=k)=\binom{k-1}{r-1}p^r(1-p)^{k-r}$
& $P(X=k)=\binom{r+k-1}{r-1}p^r(1-p)^k$
& $P(X=k)=\dfrac{\binom{w}{k}\binom{b}{n-k}}{\binom{w+b}{n}}$
& $P(X=k)=\dfrac{\binom{r+k-1}{r-1}\binom{w+b-r-k}{w-r}}{\binom{w+b}{w}}$ \\
Support
& $k=r,r+1,\dots$
& $k=0,1,2,\dots$
& $\max(0,n-b)\le k\le \min(n,w)$
& $k=0,1,\dots,b$ \\
Mean
& $\dfrac{r}{p}$
& $\dfrac{r(1-p)}{p}$
& $\dfrac{nw}{w+b}$
& $\dfrac{rb}{w+1}$ \\
Variance
& $\dfrac{r(1-p)}{p^2}$
& $\dfrac{r(1-p)}{p^2}$
& $\dfrac{w+b-n}{w+b-1}\,n\dfrac{w}{w+b}\left(1-\dfrac{w}{w+b}\right)$
& $\dfrac{rb(w+b+1)(w-r+1)}{(w+1)^2(w+2)}$ \\
MGF
& $\left(\dfrac{pe^t}{1-(1-p)e^t}\right)^r$
& $\left(\dfrac{p}{1-(1-p)e^t}\right)^r$
& $\displaystyle \sum_k e^{tk}\dfrac{\binom{w}{k}\binom{b}{n-k}}{\binom{w+b}{n}}$
& $\displaystyle \sum_{k=0}^{b} e^{tk}\dfrac{\binom{r+k-1}{r-1}\binom{w+b-r-k}{w-r}}{\binom{w+b}{w}}$ \\
\bottomrule
\end{tabularx}
\end{distbox}

\begin{center}
{\footnotesize
© 2026 Arman Jahangiri — Department of Mathematics @ University of Washington.\\
For educational use in MATH 394: Probability I
}
\end{center}

\newpage

\begin{distbox}{Continuous Distributions I}
\scriptsize
\renewcommand{\arraystretch}{1.55}
\setlength{\tabcolsep}{2.2pt}
\begin{tabularx}{\textwidth}{@{}L{2.0cm}YYYYY@{}}
\toprule
\textbf{Quantity} & \textbf{Uniform} & \textbf{Normal} & \textbf{Log-Normal} & \textbf{Exponential} & \textbf{Gamma} \\
\midrule
Parameters
& $a<b$
& $\mu\in\mathbb{R},\ \sigma^2>0$
& $\mu\in\mathbb{R},\ \sigma^2>0$
& $\lambda>0$
& $r>0,\ \lambda>0$ \\
PDF
& $f_X(x)=\dfrac{1}{b-a}$
& $f_X(x)=\dfrac{1}{\sqrt{2\pi\sigma^2}}\exp\left\{-\dfrac{(x-\mu)^2}{2\sigma^2}\right\}$
& $f_X(x)=\dfrac{1}{x\sigma\sqrt{2\pi}}\exp\left\{-\dfrac{(\log x-\mu)^2}{2\sigma^2}\right\}$
& $f_X(x)=\lambda e^{-\lambda x}$
& $f_X(x)=\dfrac{\lambda^r}{\Gamma(r)}x^{r-1}e^{-\lambda x}$ \\
Support
& $a<x<b$
& $-\infty<x<\infty$
& $x>0$
& $x\ge 0$
& $x\ge 0$ \\
Mean
& $\dfrac{a+b}{2}$
& $\mu$
& $e^{\mu+\sigma^2/2}$
& $\dfrac{1}{\lambda}$
& $\dfrac{r}{\lambda}$ \\
Variance
& $\dfrac{(b-a)^2}{12}$
& $\sigma^2$
& $e^{2\mu+\sigma^2}(e^{\sigma^2}-1)$
& $\dfrac{1}{\lambda^2}$
& $\dfrac{r}{\lambda^2}$ \\
MGF
& $\dfrac{e^{tb}-e^{ta}}{t(b-a)}$
& $\exp\left\{\mu t+\dfrac{\sigma^2t^2}{2}\right\}$
& DNE
& $\dfrac{\lambda}{\lambda-t}$
& $\left(\dfrac{\lambda}{\lambda-t}\right)^r$ \\
\bottomrule
\end{tabularx}
\end{distbox}

\vspace{2mm}

\begin{distbox}{Continuous Distributions II}
\scriptsize
\renewcommand{\arraystretch}{1.55}
\setlength{\tabcolsep}{2.2pt}
\begin{tabularx}{\textwidth}{@{}L{2.0cm}YYYY@{}}
\toprule
\textbf{Quantity} & \textbf{Beta} & \textbf{Weibull} & \textbf{Chi-Square} & \textbf{Student-$t$} \\
\midrule
Parameters
& $a,b>0$
& $\alpha,\lambda>0$
& $n>0$
& $n>0$ \\
PDF
& $f_X(x)=\dfrac{\Gamma(a+b)}{\Gamma(a)\Gamma(b)}x^{a-1}(1-x)^{b-1}$
& $f_X(x)=\alpha\lambda x^{\alpha-1}e^{-\lambda x^\alpha}$
& $f_X(x)=\dfrac{1}{2^{n/2}\Gamma(n/2)}x^{n/2-1}e^{-x/2}$
& $f_X(x)=\dfrac{\Gamma((n+1)/2)}{\sqrt{n\pi}\Gamma(n/2)}\left(1+\dfrac{x^2}{n}\right)^{-(n+1)/2}$ \\
Support
& $0<x<1$
& $x>0$
& $x>0$
& $-\infty<x<\infty$ \\
Mean
& $\dfrac{a}{a+b}$
& $\dfrac{\Gamma(1+1/\alpha)}{\lambda^{1/\alpha}}$
& $n$
& $0,\ n>1$ \\
Variance
& $\dfrac{ab}{(a+b)^2(a+b+1)}$
& $\dfrac{\Gamma(1+2/\alpha)}{\lambda^{2/\alpha}}-\mu^2$
& $2n$
& $\dfrac{n}{n-2},\ n>2$ \\
MGF
& --
& --
& $(1-2t)^{-n/2}$
& DNE \\
\bottomrule
\end{tabularx}
\end{distbox}

\begin{distbox}{Moment-Generating Function}
The moment-generating function is
\[
M_X(t)=\mathbb{E}(e^{tX}).
\]
When it exists near $t=0$, it can be used to compute moments:
\[
\mathbb{E}(X^n)=M_X^{(n)}(0).
\]
\end{distbox}

\begin{center}
{\footnotesize
© 2026 Arman Jahangiri — Department of Mathematics @ University of Washington.\\
For educational use in MATH 394: Probability I
}
\end{center}

\end{document}
